\documentclass{article} % For LaTeX2e
\usepackage{iclr2026_conference,times}

% Optional math commands from https://github.com/goodfeli/dlbook_notation.
\input{math_commands.tex}

\usepackage{hyperref}
\usepackage{url}

\title{Origami as a Spatial Reasoning Benchmark}

% Authors must not appear in the submitted version. They should be hidden
% as long as the \iclrfinalcopy macro remains commented out below.
% Non-anonymous submissions will be rejected without review.
%
% With \iclrfinalcopy commented out, iclr2026_conference.sty replaces this
% block with "Anonymous authors / Paper under double-blind review", so the
% submission is anonymous regardless of what is written here.
\author{Anonymous authors \\
Paper under double-blind review}

\newcommand{\fix}{\marginpar{FIX}}
\newcommand{\new}{\marginpar{NEW}}

%\iclrfinalcopy % Uncomment for camera-ready version, but NOT for submission.
\begin{document}


\maketitle

\begin{abstract}
Progress on reasoning has been fastest where verification is free. In mathematics and in code, a
candidate answer can be checked exactly, cheaply, and without a human, and that property has
shaped how quickly those domains improved. Spatial and physical reasoning has no such checker, so
it is evaluated instead with static images, multiple choice questions, or learned surrogate
simulators whose wrong verdicts look exactly like their right ones. A benchmark can only be as
trustworthy as the judge inside it. We show that origami is a domain where exact verification is
free, and that the gap between generating a problem and solving it is severe. Folding a sheet
forward is trivial: fold, record, unfold. Recovering the discrete sequence of folds that produced
a crease pattern is NP hard. The same forward engine that generates a sample can therefore rule on
any proposed step exactly, at no cost, with no learned model in the loop. The task is also not a
puzzle chosen for convenience. It becomes strictly harder as it proceeds, because every fold
thickens the stack and constrains the next one, and a solver has to track a planar geometry
together with a layer ordering that continuously constrain each other. We release CP2Seq, a
benchmark whose ground truth is constructed rather than annotated, together with an environment
that is also the verifier: stepping it is the same act as asking for a verdict. It executes one
candidate fold and either returns the resulting state or refuses with a named reason, such as the
sheet would tear, the fold creases nothing, or the selected layers cannot move in that direction.
It performs no search and makes no choices of its own, so the work of proposing, pruning and
backtracking stays with the model, and what the benchmark scores is search rather than geometry.
Samples are generated under two action models taken from the simple folding literature and
stratified along two axes, fold depth and coupling, where coupling is the number of layers a
single fold cuts. Every sample rebuilds byte identically from its seed and is checked by tolerance
free replay. Because one crease pattern admits many valid folded states, and because the recorded
sequence is not guaranteed to be minimal, an answer is scored by replaying the proposed sequence
rather than by step wise agreement with the stored one. The scoring protocol is solve rate at a
fixed query budget, reported per difficulty stratum under two notions of equality, equality of
crease sets and equality of folded states up to a symmetry group declared in advance, with queries
to solution as a secondary measure over all attempts including those that exhaust the budget. We
instantiate this protocol on off the shelf multimodal language models, with ablations that remove
visual feedback, remove filtering, and remove tools entirely as a memorization control on
anonymized geometry.
\end{abstract}

\section{Positioning}
\label{sec:positioning}

This paper is submitted as a \textbf{benchmark and dataset} contribution. What it offers the
community is an evaluation substrate for \textbf{spatial and geometric reasoning in large language
models}: a corpus, an environment that doubles as an exact verifier, and a scoring protocol,
rather than a model, a training method or a state-of-the-art number. Every arm reported here runs
an off-the-shelf model through the harness; nothing is trained. The claim under review is that the
benchmark measures what it says it measures and that its judge can be trusted, not that any
particular model is good at the task.

The reasoning being measured is geometric as much as it is spatial, and the distinction is
load-bearing rather than cosmetic. A solver has to reason about lines, reflections, incidence and
adjacency on the original sheet, which is geometry in the ordinary sense, and simultaneously about
a layer ordering that the geometry does not determine, which is combinatorial. Calling the task
spatial reasoning alone understates the second half. Positioning the work as a benchmark for
spatial \emph{and geometric} reasoning states the scope accurately and places it beside the
multimodal-reasoning benchmarks of Section~\ref{sec:related-benchmarks} rather than beside the
computational-origami literature of Section~\ref{sec:related-theory}, which is the theory the
benchmark rests on rather than the field it contributes to.

% TITLE, to be settled before submission. The title may carry "geometric reasoning" explicitly.
% Candidates, mirrored from the Positioning section of paper/DRAFT.md:
%   - Origami as a Spatial and Geometric Reasoning Benchmark
%   - Origami as a Benchmark for Spatial and Geometric Reasoning in Multimodal Language Models
%   - CP2Seq: A Spatial and Geometric Reasoning Benchmark with an Exact Verifier
% The subtitle "with Free Exact Verification" was removed from the title; the exactness of the
% verifier now carries in the abstract and Section 5 rather than in the title.

\section{Introduction}
\label{sec:intro}

Reasoning has improved fastest in the two domains where an answer can be checked without a human.
A proof can be run through a proof assistant; a program can be run against its tests. In both
cases the check is exact, it is cheap enough to run millions of times, and it does not itself have
to be trusted, because it is not a model of the thing, it is the thing. That property is what
makes large scale evaluation, rejection sampling, search at test time and reward from verification
possible at all.

Spatial and physical reasoning has no such checker, and the field has adapted in three ways, each
of which costs something specific. Benchmarks built on static images or multiple choice questions
never execute the model's proposal, so an answer that is right for the wrong reason is
indistinguishable from one that is right. Benchmarks and methods built on learned surrogate
simulators do execute the proposal, but against an approximation: a wrong verdict looks exactly
like a right one, which means the benchmark's own error rate is unknown and unmeasurable from
inside. Human judgement and preference scores are exact in a sense, but they do not scale and they
are not reproducible across papers. The common consequence is worth stating plainly, because it is
the motivation for everything below: \textbf{a benchmark can only be as trustworthy as the judge
inside it.}

We observe that origami has the property the field is missing, and that it has it in an unusually
clean form. Folding a sheet forward is trivial: pick a line, fold across it, record what you did,
and unfold at the end to read off the crease pattern. Recovering the discrete sequence of folds
that produced a given crease pattern is NP hard \citep{arkin-map,akitaya-hard}. The forward
direction is not merely easier than the inverse; it is the \emph{generator}, and the same engine
that generates a problem can rule on any proposed step of a solution exactly. Verification costs
one forward application of an operation the engine already implements. There is no learned
component anywhere in the loop, no tolerance to tune, and no oracle to trust.

This asymmetry is what a benchmark wants. Difficulty is controllable rather than sampled and hoped
for, because the generator chooses how deep and how tangled a sample is. Ground truth is free
rather than annotated, because the sequence is recorded as it happens rather than inferred
afterwards. Negative examples are unlimited, because any illegal fold is generated by asking for
one. And the judge is the same object as the environment, so the benchmark's verdicts do not
depend on a second artifact that could disagree with the first.

The task is also hard in ways that a maze or a grid-world is not, and the differences are
measurable rather than rhetorical. First, it gets harder as it proceeds: every fold thickens the
stack and constrains what the next fold may do, so the difficulty of step $k$ is a function of
what happened in steps $1$ through $k-1$. Second, a solver has to maintain two kinds of state that
constrain each other. The crease graph is fixed once the pattern is given, but the vertex
coordinates change at every step, the overlap structure changes at every step, and the layer
ordering changes at every step and is \emph{chosen} rather than computed. Determining that
ordering is NP hard even when a valid mountain-valley assignment is supplied \citep{bern-hayes}.
Mazes have geometry alone; here an independent combinatorial layer has to be tracked alongside it.
Third, the action space is not something we invented for the occasion: simple folds have been
defined and their complexity studied for two decades \citep{arkin-map}, so the operations an agent
is given have names, and claims about them can be checked against a literature.

We contribute the following.

\begin{enumerate}
\item \textbf{CP2Seq}, a benchmark for crease-pattern-to-sequence recovery whose ground truth is
  constructed rather than annotated, generated under two named action models from the simple
  folding literature, stratified by fold depth and by coupling, and reproducible from seeds
  rather than shipped as data.
\item \textbf{An environment that is also the verifier.} Stepping it is the same act as asking for
  a verdict. It is exact, it is free, it performs no search, and its refusals are named rather
  than boolean, so a rejection is a diagnosis instead of a dead end.
\item \textbf{A scoring protocol that survives non-unique solutions.} Answers are judged by
  replay, at two levels of equality, with the symmetry group fixed in advance, because the
  recorded sequence is one solution among several and is not guaranteed to be the shortest.
\item \textbf{An evaluation of off-the-shelf multimodal language models} in this loop, with a
  visual-feedback ablation, a filtering ablation, and a memorization control on anonymized
  geometry. \texttt{[TO RUN]}
\item \textbf{A measured scope boundary.} 89.3\%% fact:probeC.provenNotPct
  of real crease patterns provably lie outside all-layers simple folding. The benchmark's action
  space is bounded by evidence rather than by assertion, and the boundary is reported rather than
  buried.
\end{enumerate}

\section{Related work}
\label{sec:related}

\subsection{Origami benchmarks for multimodal models}
\label{sec:related-benchmarks}

\textbf{OrigamiSpace} \citep{origamispace} contributes 350 instances, each carrying a crease
pattern diagram, a compiled flat pattern, the complete folding process and a final folded image,
and defines four tasks over them: pattern prediction, multi-step spatial reasoning, spatial
relationship prediction, and end-to-end crease-pattern code generation. It also provides an
interactive environment and describes reinforcement learning as a possibility explored rather than
a result demonstrated. Its design lesson is that one richly instrumented sample can support
several tasks, which is how 350 instances sustain a benchmark paper.

\textbf{GamiBench} \citep{gamibench} contributes 186 regular and 186 \emph{impossible} crease
patterns paired with folded shapes from six viewpoints, across three visual question-answering
tasks, and introduces two diagnostic metrics: viewpoint consistency and impossible fold selection
rate. Its impossible patterns are the closest the current literature comes to forcing a model to
judge feasibility rather than recognise a shape, and its headline finding is that leading models
struggle even at single-step spatial understanding. It is the benchmark this work is most directly
in conversation with, and the one whose negative examples we can produce without limit, because
generating an illegal fold in our environment is the same operation as generating a legal one.

\textbf{OrigamiBench} \citep{origamibench}, \textbf{FoldingAgent} \citep{foldingagent} and
\textbf{COrigami} \citep{corigami} complete the picture; their metrics appear in the comparison of
Section~\ref{sec:related-gap}.

What all of these share is the gap. In each, the model's proposal is answered by a question rather
than by an execution: it selects among rendered alternatives, or produces an artifact that is
compared to a target. None of them steps a paper-exact engine and asks whether \emph{this} move,
on \emph{this} stack, is something paper could do.

\subsection{Methods that pair a proposer with a simulator}
\label{sec:related-proposer}

\textbf{Learn2Fold} \citep{learn2fold} formulates origami as conditional program induction over a
crease-pattern graph. A language model generates candidate folding programs from text, and a
learned graph-structured world model acts as a differentiable surrogate simulator that predicts
physical feasibility and failure modes before execution, inside a lookahead planning loop. Its
stated key insight is to decouple semantic proposal from physical verification.

We adopt that decomposition and differ on one thing: the verifier. Theirs is learned,
differentiable and approximate, which is what makes it usable for planning and also what makes its
mistakes invisible. Ours is exact and costs nothing, because it is the generator run forwards.
Their scoring is step-level precision, recall and F1 against an expert sequence, together with
edge-level IoU; Section~\ref{sec:scoring-why} shows why that family of measures is unsafe in this
task, since a model that finds a \emph{shorter} correct sequence is marked wrong by it. That is a
finding about scoring in this setting rather than a criticism of their results.

The general precedent for a proposer paired with a verifier is older than any of this work: a
network that proposes and a search that checks is the structure behind AlphaGo, and the reason the
combination is stronger than either half.

\subsection{Computational origami: the theory under the action space}
\label{sec:related-theory}

The benchmark's central claim, that geometry does not determine a folded state, is Bern and
Hayes's: deciding the overlap order of a flat folding is NP hard even given a valid
mountain-valley assignment \citep{bern-hayes}. The action space is Arkin et al.'s: simple folds,
in the one-layer, some-layers and all-layers variants, with finite or infinite fold lines
\citep{arkin-map}. Their complexity has been mapped since, including simple foldability's hardness
\citep{akitaya-hard} and the infinite all-layers case \citep{akitaya-infinite}. The solver whose
constraint vocabulary we borrow for terminal states is Flat-Folder \citep{akitaya-flatfolder}.

Two results shape what the task actually is. \citet{continuous} showed that any well-behaved
folded state is reachable by \emph{some} continuous motion; the question is therefore never
whether a path exists, but whether one exists through a finite number of the discrete operations a
hand or a machine can perform. And counting foldings has no closed form even for a
one-dimensional strip, a problem open since Lucas attributed it to Lemoine in 1891
\citep{lucas,oeis-a000136}; this is the answer to ``why not enumerate the states'', and it is a
fact about the problem rather than about any implementation.

Finally, the problem itself is not new: generating folding sequences from crease patterns was
named and attacked symbolically by \citet{akitaya-cp2seq}, by the same group whose tools this
field now depends on. What is new is neither the problem nor the theory, but the observation that
this problem supplies an exact verifier for free and can therefore serve as an evaluation
substrate for models that reason with tools.

\subsection{The structural gap}
\label{sec:related-gap}

Collecting the metrics used across this literature makes one absence visible. Every metric in use
compares a produced artifact against a target: geometric or semantic similarity to a reference,
precision and recall against an expert sequence, IoU over affected edges, compilation success,
human preference. \textbf{None ranks among the several valid folded states of a single crease
pattern, and none scores an agent by how much verification it required.} The absence is structural
rather than accidental, because computing any of those metrics requires a target to compare
against, and the interesting question here has more than one right answer.

\section{Problem formulation}
\label{sec:problem}

\subsection{Flat-folded states}
\label{sec:problem-states}

A flat-folded state is an isometric map of the sheet into the plane together with a layer
ordering. The map is fixed by the crease pattern; the ordering is not. A single crease pattern
therefore admits many folded states that are geometrically identical and distinct as embeddings,
and determining the ordering is NP hard even when a valid mountain-valley assignment is given
\citep{bern-hayes}.

This distinction is the reason the task cannot be reduced to geometry. Over a folding sequence the
crease graph's topology never changes: the same vertices, edges and faces exist at step one and at
step nineteen. What changes is everything else. Vertex coordinates are reflected at every step.
The overlap structure, which faces lie above which, changes at every step and follows from the
geometry. The layer ordering changes at every step and does \emph{not} follow from the geometry;
it is chosen, and it is where the decisions live. The number of layers accumulates, which is the
formal content of ``it gets harder as it goes''.

\subsection{The action space}
\label{sec:problem-actions}

Simple folds come in named variants, and any benchmark in this area has to say which it uses,
because the complexity results differ between them.

\begin{table}[h]
\caption{Named variants of the simple fold \citep{arkin-map}.}
\label{tab:action-space}
\begin{center}
\begin{tabular}{ll}
\multicolumn{1}{c}{\bf MODEL}  &\multicolumn{1}{c}{\bf DEFINITION} \\
\hline \\
One-layer simple fold    & Rotate a single layer $\pm 180^\circ$ about a line \\
Some-layers simple fold  & Rotate a contiguous run of layers \\
All-layers simple fold   & Rotate every layer the line crosses, as sheet metal bends \\
Finite vs infinite line  & Whether the fold line is a full line or a segment \\
\end{tabular}
\end{center}
\end{table}

CP2Seq instantiates the all-layers and some-layers tiers. The all-layers tier is the simpler
object and it buys three properties at once, the third of which is rarely named. Because the whole
stack moves as a rigid body, no two connected pieces of paper ever move relative to each other
except at the fold line itself, where paper is allowed to bend. Self-intersection is impossible,
the layer ordering is forced rather than chosen, and \textbf{tearing is impossible by
construction.}

Moving only some layers brings all three back, and tearing becomes the binding constraint. If a
moving face is joined to a stationary face along a crease, and that crease is not on the fold
line, the fold would rip the sheet. This is decidable only in the coordinates of the original
square, because adjacency is a fact about the original sheet that survives any amount of folding
and cannot be recovered from the current positions of the polygons alone.

An example makes the tier difference concrete, and it needs only two folds. Fold the right half of
a square over to the left across the vertical midline. The result is a two-layer stack whose top
layer is joined to the bottom layer along the entire crease at the midline. Now select the top
layer alone. Folding its upper half down across the horizontal midline would tear the sheet,
because the moving piece is joined to stationary paper along a vertical crease segment that is
perpendicular to the fold line. Folding its left quarter across a vertical line does not tear,
because the moving piece's only join to stationary paper lies on the fold line itself. Same state,
same layer selection, opposite verdicts, and the only difference is the orientation of the line.

\subsection{The task}
\label{sec:problem-task}

Given a crease pattern, produce a sequence of simple folds that reproduces it. Generating an
instance costs one forward pass of the engine; solving it is NP hard. That asymmetry is the
paper's foundation.

\textbf{Caveat.} One claim must not be made, and it was drafted wrongly here once. It is tempting
to say that legal moves are rare and that this is what makes the task hard. Enumeration refutes
it: legal partial folds \emph{grow} with depth, from 22 at two layers to 332 at thirty-eight, six
times the 56 all-layers folds available at the same state. What falls is the hit rate of uniform
random proposal, from 11.9\% to 3.8\% by seventeen layers, because the space being sampled grows
faster than the legal set inside it. That is a fact about a sampler, not about origami, and
conflating the two would put a false statement about branching factor into the paper.

\section{CP2Seq}
\label{sec:cp2seq}

\subsection{Generation}
\label{sec:cp2seq-generation}

The corpus is synthesised and the project takes no outside data. A sample is produced by folding
forward from a square with randomly chosen simple folds, recording each fold as it is applied, and
unfolding at the end to obtain the crease pattern. The \texttt{(pattern, sequence)} pair is built
rather than labelled, so there is nothing to annotate and nothing to trust, and every sample is
correct by construction: the pattern is what the folding left behind.

Nothing is shipped that cannot be rebuilt. Given its seed, a sample regenerates byte-identically,
so the corpus is distributed as a generator plus a manifest rather than as data.

\subsection{Difficulty, stratified rather than sampled}
\label{sec:cp2seq-difficulty}

Difficulty is set by the sampler along two axes rather than measured after the fact. \textbf{Fold
depth} is the primary axis. \textbf{Coupling}, the number of creases a single fold creates, which
is to say how many layers it cuts, is the second; it is free to compute at generation time, since
one fold cuts every layer it crosses and each cut becomes one crease in the unfolded square.
Coupling is simultaneously the non-local dependency that Learn2Fold's difficulty tiers appeal to
and the closest measurable proxy for how hard a model is to fold by hand.

The release corpus holds 600 samples in five batches: 400 all-layers samples spanning 4 to 19
folds, stratified easy, mid and hard on the action space's tertiles, and 200 some-layers samples
at depths 3 to 6. Median crease count is 33 and median layer count 24, with the deep tail reaching
tens of thousands of both.

\textbf{Caveat.} Two limits belong in the paper rather than in an appendix. The first is an empty
cell: long sequences at low coupling are almost unreachable, yielding 2% fact:corpus.synth.lLocalHits
hits in 16,000 attempts, because every all-layers fold thickens the stack, so a long sequence
cannot keep cutting few layers. Real folders reach that corner by \textbf{pre-creasing}, an
operation this action space excludes. The longer a real model runs, the more of it sits outside
what the benchmark can express. The second is depth: the some-layers tier stops at six folds in
this release, and the difficulty the benchmark is about lives past the depth a search can reach.
Restoring the deeper tier is a matter of generator configuration and wall-clock, not redesign.

\subsection{What a sample carries}
\label{sec:cp2seq-sample}

\begin{table}[h]
\caption{The files shipped with each CP2Seq sample.}
\label{tab:sample-files}
\begin{center}
\begin{tabular}{ll}
\multicolumn{1}{c}{\bf FILE}  &\multicolumn{1}{c}{\bf CONTENTS} \\
\hline \\
\texttt{cp.fold}    & The crease pattern, planarised. The input \\
\texttt{seq.json}   & Every fold: line, direction, creases made, coupling. The ground truth \\
\texttt{steps.fold} & The pattern plus the folded state after each step, as one multi-frame file \\
\texttt{meta.json}  & Difficulty metrics, degeneracy flags, provenance \\
\end{tabular}
\end{center}
\end{table}

\subsection{Verifying the corpus itself}
\label{sec:cp2seq-verify}

Two checks run over every sample, written independently of each other. The first replays each
recorded sequence and confirms it reproduces the pattern stored beside it. The second, written
without reference to the first, compares maximal creased intervals as multisets with no tolerance
in the verdict, after clustering edges into lines so that subdivision differences cannot register
as disagreements. Both pass on all 600 samples of the release corpus.

\textbf{Caveat.} Both share the fold engine with the generator, which bounds what they can
establish. They cannot catch a wrong model of paper: if the engine is wrong about tearing, replay
is wrong in the same way and agrees with itself. What they do catch is drift between what was
folded and what was recorded, which is the likelier failure and the one that silently corrupts
ground truth. An independent check requires a third-party solver and is not claimed here.

\section{The environment, which is also the verifier}
\label{sec:env}

\subsection{One object, two names}
\label{sec:env-one-object}

In this benchmark the environment and the verifier are the same object. Stepping it is the same
act as asking for a verdict: the model submits a candidate fold, and the reply is either the
resulting state or a refusal. There is no separate validity oracle, no second artifact that could
disagree with the first, and no learned component in the loop.

\subsection{The division of labour}
\label{sec:env-division}

\begin{table}[h]
\caption{Who does what. The environment does not search.}
\label{tab:division}
\begin{center}
\begin{tabular}{ll}
\multicolumn{1}{c}{\bf COMPONENT}  &\multicolumn{1}{c}{\bf RESPONSIBILITY} \\
\hline \\
Environment / verifier & State update and legality. It does not search and does not choose \\
Model                  & Proposes the next fold, decides where to go, recognises when to backtrack \\
\end{tabular}
\end{center}
\end{table}

This division is what makes the benchmark measure something. The obvious objection to any
tool-augmented setup is that the tool is doing the work, and here the answer is structural rather
than rhetorical: deciding global flat-foldability is NP hard, so the environment cannot search its
way to an answer even in principle. It can only say whether one proposed step is something paper
could do. Compressing an exponential search into a feasible number of queries is exactly what the
environment cannot supply and exactly what the model is being measured on. The referee is not a
player.

\subsection{Refusals are named}
\label{sec:env-refusals}

An illegal fold returns a structured refusal with a name, not a boolean.

\begin{table}[h]
\caption{The refusal taxonomy. A rejection arrives already labelled with the class it violated.}
\label{tab:refusals}
\begin{center}
\begin{tabular}{lp{0.66\textwidth}}
\multicolumn{1}{c}{\bf REFUSAL}  &\multicolumn{1}{c}{\bf MEANING} \\
\hline \\
\texttt{would-tear} & The moving layers are joined to stationary paper away from the fold line, so
  the sheet would come apart. The binding constraint of the some-layers tier \\
\texttt{no-crease} & The fold moves layers without creasing anything, tearing them free of the
  sheet rather than folding them \\
\texttt{nothing-to-move} & The fold line misses the selected layers entirely \\
\texttt{direction-impossible} & A run taken from the bottom of the stack was asked to fold over,
  or the reverse \\
\texttt{bad-line} & The fold line is malformed, or names no side of itself \\
\end{tabular}
\end{center}
\end{table}

The names are the design. A model told \texttt{would-tear} can act on the information; a model told
``illegal'' cannot. This also turns error analysis into a taxonomy that exists before the
experiment: every rejection arrives already labelled with the class it violated, so a reject is a
diagnosis rather than a dead end.

\textbf{Caveat.} There is deliberately no self-intersection refusal, and there cannot be one in
this action space. A some-layers fold may only move a contiguous run of layers taken from the top
or the bottom of the stack, and such a move cannot drive paper through the layers it left behind.
Lifting the top two layers and folding them across is something a hand can do; folding them
\emph{underneath} the stack is not a fold but a slit. Restricting selection to a run at one
extremity makes the illegal case unrepresentable rather than undetected. The four constraint
classes that Flat-Folder checks (taco-taco, taco-tortilla, tortilla-tortilla, transitivity)
\citep{akitaya-flatfolder,flatfolder} remain the right vocabulary for global terminal states; they
are not what a single legal step needs checking against here.

\subsection{What the model observes}
\label{sec:env-observes}

On success the environment returns the new state together with two views: a top-down X-ray of the
stack, and an exploded view of the layers. Every intermediate state in this tier is flat, so a
second camera angle would show the same silhouette rotated and would carry no new information;
what carries information is layer structure, which is why the second view is an explosion rather
than a rotation. Positions are reported as coordinates on the original sheet together with an
integer layer index. There is no continuous vertical coordinate at this tier, and a rendering that
looks three-dimensional is not evidence of one; thickness belongs to a different problem and a
different paper.

\subsection{Error verbosity is an experimental variable}
\label{sec:env-verbosity}

The richer a refusal, the more of the reasoning the environment performs rather than the model. At
the limit, a message that names the fix has solved the step. Verbosity is therefore fixed before
the first run and reported with the results; otherwise the tool-augmented condition is not
reproducible and its comparison against the verifier-only condition measures message design rather
than reasoning.

\section{Scoring}
\label{sec:scoring}

\subsection{Why the obvious protocol is wrong}
\label{sec:scoring-why}

The recorded sequence is \emph{a} solution, not \emph{the} solution. The generator folds; it does
not search, so nothing makes its sequence minimal, and no minimal-length label is shipped because
producing one would require a search this benchmark deliberately does not run. Measured on the
verified split, 2% fact:shorter.release.n
of 100% fact:shorter.release.outOf
samples admit a sequence one fold shorter than the one that built them.

The consequence is sharp. \textbf{Any metric that compares a proposal step by step against the
recorded sequence marks a better answer wrong.} A model that finds a shorter correct solution is
penalised for it. This disqualifies edit distance, and it disqualifies step-level precision,
recall and F1 against an expert sequence for this task. Step count is not a correctness signal in
either direction.

\subsection{What we do instead}
\label{sec:scoring-what}

Proposals are scored by replaying them through the engine and comparing the resulting states. Both
sides are replayed: the model's sequence and the reference sequence both pass through the same
engine, and a stored frame is never the object compared against. That last point is a consequence
of the file format rather than a preference. Stored frames record where each face sits in the
current plane and not which piece of the original sheet it is, and deep in a stack many layers are
congruent triangles, so two states differing by a swap of two such layers cannot be distinguished
from stored frames alone. A replayed state carries original-sheet coordinates for free, because
the engine needs them to decide tearing at all.

Two levels of equality are reported.

\textbf{Level 1, crease sets.} An exact multiset comparison of maximal creased intervals, equal
rather than overlapping, with no tolerance in the verdict.

\textbf{Level 2, folded states up to a declared group.} Because one crease pattern admits many
valid terminal states, exact equality against the one stored state would mark correct answers
wrong. The answer is not a tolerance but an equivalence: declare in advance which differences do
not count, then demand exactness inside that. The group is the eight symmetries of the square
together with an arbitrary translation. Translation is included because a fold can carry paper off
the original square, so a folded state's position in the plane is an accident of the sequence. The
four reflections additionally reverse the stack and flip every face's parity, because turning a
model over does all three at once; applying a reflection to coordinates alone would compare a
model against its mirror image and silently accept wrong answers.

\textbf{Caveat.} The layer-order variants that a flat-foldability solver reports as equally valid
are deliberately \emph{not} in the group. Those are not symmetries of one state; they are
different states reachable from the same crease pattern, and deciding which orderings are valid
requires a solver. Folding that into an equality test would hide an external dependency inside
what reads as arithmetic. If that equivalence is wanted, it belongs in a separate check that owns
the solver.

\subsection{Metrics}
\label{sec:scoring-metrics}

The primary metric is \textbf{solve rate at a fixed query budget}, reported per difficulty stratum
and never pooled, under both levels of equality. A proposal is solved if replaying it reproduces
the target. Exhausting the budget is a timeout, which is reported separately from a wrong answer
because the two mean different things.

The secondary metric is \textbf{queries to solution}, reported over all attempts including those
that time out. Conditioning on the solved attempts hides precisely the long-tail blowup that makes
the task interesting. This metric is what answers the ``the simulator did all the work'' objection
quantitatively: two models sharing one exact verifier can be ranked by how little of it they
needed.

\subsection{Validating the protocol}
\label{sec:scoring-validate}

A verifier paper that does not test its own verifier is asking to be taken on trust, so the
comparator ships with a negative control: a deliberately corrupted state, with two layers swapped,
must be rejected.

The control earned its place. On its first full run, 23 of 600 corrupted states were
\textbf{accepted}. The cause was not a tolerance. The comparison used current-plane geometry and
parity alone, and deep in a stack many layers are congruent triangles, so two layers occupying the
same region but originating from different parts of the sheet were interchangeable. They are not
interchangeable: which one lies underneath is a real difference between two folded states, and
precisely the difference Level 2 exists to catch. With face identity on the original sheet
restored to the comparison, the full corpus reports 600 states equal, all under the identity
symmetry, and 600 of 600 corrupted states rejected. Every verdict additionally reports whether
face identity was available, so a weakened comparison is visible in the output rather than assumed
away.

\section{Objections, answered here rather than in rebuttal}
\label{sec:objections}

\textbf{``The simulator does all the work.''} It cannot: deciding flat-foldability is NP hard, and
the environment only rules on single steps. Section~\ref{sec:env-division} states the division of
labour and Section~\ref{sec:scoring-metrics} measures it.

\textbf{``Origami is a toy problem.''} Flat origami is Turing complete \citep{turing}, so the
domain is not expressively impoverished. \textbf{Caveat.} This is a defensive citation and nothing
more. It shows that the domain is rich enough to be interesting; it does not show that anything
learned here transfers, and the causal chain from ``origami is Turing complete'' to ``training on
origami yields physical understanding'' is broken in the middle. Rule 110 is Turing complete too.

\textbf{``Why not enumerate the valid states?''} There is no closed form for the number of foldings
even of a one-dimensional strip, a problem open since 1891 \citep{lucas,oeis-a000136,koehler}, and
state counts for real patterns reach astronomical magnitudes. The hardness is the problem's, not
an implementation's.

\textbf{``Synthetic data is a shortcut.''} The comparable published corpora are largely produced by
their authors' own symbolic simulators; synthesis is the field's normal practice. Here it is also
what makes difficulty controllable and ground truth free, and what allows the corpus to ship as a
seed rather than as a file.

\textbf{``Your corpus is not real origami.''} Correct, and measured: 89.3\%% fact:probeC.provenNotPct
of real crease patterns provably lie outside all-layers simple folding. This is stated as a scope
boundary in Section~\ref{sec:intro} and Section~\ref{sec:limitations} rather than left for a
reviewer to discover.

\textbf{``The model may have memorised the pattern.''} The no-tools arm runs on anonymized geometry
with filenames and metadata stripped, for exactly this reason: if a model can read a pattern's name
it recalls rather than reasons, and then both arms measure recall and the comparison measures
nothing. See Section~\ref{sec:ablations}.

\section{Experiments}
\label{sec:experiments}

\texttt{[TO RUN]}

\begin{itemize}
\item \textbf{Models.} Held fixed across arms; comparing across models would confound the
  ablation. \texttt{[TO RUN: the main evaluation sweep]}
\item \textbf{Protocol.} Query budget, repeats, seeds; median and spread reported, never a single
  run, because sampling is not deterministic and one run is not a measurement.
  \texttt{[TO RUN: same sweep]}
\item \textbf{Main table.} Solve rate per stratum under Level 1 and Level 2.
  \texttt{[TO RUN: same sweep]}
\item \textbf{Query distribution.} Queries to solution over all attempts, with the timeout
  fraction. \texttt{[TO RUN: same sweep]}
\end{itemize}

\section{Ablations}
\label{sec:ablations}

\texttt{[TO RUN]}

\begin{table}[h]
\caption{The ablation arms. Same model, same patterns, same query budget across arms.}
\label{tab:ablations}
\begin{center}
\begin{tabular}{llp{0.34\textwidth}}
\multicolumn{1}{c}{\bf ARM}  &\multicolumn{1}{c}{\bf TOOLS AVAILABLE}  &\multicolumn{1}{c}{\bf WHAT IT MEASURES} \\
\hline \\
Full tool belt & Environment, views, any auxiliary tools & Upper bound \\
No vision      & Same, rendered views withheld           & The value of visual feedback specifically \\
Verifier only  & Pass/fail, no filtering                 & The value of filtering on top of raw verification \\
No tools       & None, anonymized geometry               & Memorization control \\
\end{tabular}
\end{center}
\end{table}

Same model, same patterns, same query budget across arms. If the no-vision or no-tools arm
performs nearly as well as the full arm, \textbf{that is a finding rather than a failed
experiment}: it means the gain is not where it was assumed to be. \texttt{[TO RUN]}

\section{Results}
\label{sec:results}

\texttt{[TO RUN]}

\begin{itemize}
\item \textbf{Headline.} \texttt{[TO RUN]}
\item \textbf{Failure taxonomy.} Rejections grouped by refusal class; the taxonomy already exists,
  so this is tabulation rather than interpretation. \texttt{[TO RUN]}
\item \textbf{Where difficulty lives.} Depth against coupling. \texttt{[TO RUN]}
\end{itemize}

\section{Limitations}
\label{sec:limitations}

\textbf{No training.} Every arm is an off-the-shelf model driven by prompting and tool calling; no
model is trained or fine-tuned. The contribution is the harness, not a model. This is stated in the
introduction as well as here, because a limitation the authors declare is context and one a
reviewer discovers is an objection.

\textbf{Scope of the action space.} 89.3\%% fact:probeC.provenNotPct
of real crease patterns lie outside all-layers simple folding, and pre-creasing, the operation real
folders use to reach long sequences at low coupling, is excluded.

\textbf{The corpus verifiers share an engine with the generator}, so they cannot detect a wrong
model of paper. They detect drift between what was folded and what was recorded.

\textbf{Flat states only.} No thickness, no material, no mechanics. Thickness changes the geometry
outright and belongs to a separate line of work with a verifier several orders of magnitude more
expensive.

\textbf{Level 2's group is narrower than the truest notion of equality}, because valid layer-order
variants are excluded for the reason given in Section~\ref{sec:scoring-what}. It is therefore
strict rather than permissive, which is the safe direction but not a free one.

\section{Figures}
\label{sec:figures}

\begin{table}[h]
\caption{Planned figures and their status.}
\label{tab:figures}
\begin{center}
\begin{tabular}{llp{0.5\textwidth}l}
\multicolumn{1}{c}{\bf \#}  &\multicolumn{1}{c}{\bf FIGURE}  &\multicolumn{1}{c}{\bf STATUS} \\
\hline \\
1 & The asymmetry: generation is one forward pass, recovery is NP hard & To draw \\
2 & The tearing pair of Section~\ref{sec:problem-actions}. One state, one layer selection, a
    vertical fold line legal and a horizontal one not & To draw \\
3 & A refusal as the model receives it: view, named reason, per-constraint checklist with
    locations & To draw \\
4 & The difficulty grid, depth against coupling, with the empty cell visible & To draw \\
5 & Main result per stratum & \texttt{[TO RUN]} \\
\end{tabular}
\end{center}
\end{table}

\section{Reproducibility}
\label{sec:reproducibility}

The corpus ships as a generator and a manifest; every sample rebuilds byte-identically from its
seed. The environment, both corpus verifiers and the state comparator are released with the
benchmark. The negative control of Section~\ref{sec:scoring-validate} runs as a command, so a
reader can confirm that the comparator rejects corrupted stacks rather than taking
Section~\ref{sec:scoring-validate} on trust.

\subsubsection*{Acknowledgments}

Withheld for anonymous review.

\bibliography{iclr2026_conference}
\bibliographystyle{iclr2026_conference}

\appendix

\section{Reference notes}
\label{app:references}

The bibliography is grouped in \texttt{iclr2026\_conference.bib} by the role each work plays in
the argument rather than alphabetically, so that a co-author can see what is load-bearing. The
notes below record what is still outstanding.

\paragraph{Origami benchmarks and systems for language models.}
\citet{origamispace} must be cited as a preprint; \texttt{papers.md} recorded NeurIPS'25 and the
arXiv record carries no venue. \citet{gamibench} has code and data released. \citet{foldingagent}
reconstructs folding processes from instructional video and is built on the PurelandFold dataset;
code at \url{https://github.com/maya-moriya/FoldingAgent}. \textbf{OrigamiBench}
\citep{origamibench} has a \emph{citation missing}: it appears in our metrics survey as the source
of query efficiency and geometric/semantic similarity, and its dataset is the
366% fact:corpus.instagram.total
patterns in Flat-Folder's \texttt{examples/instagram/}. A full reference has to be recovered before
submission. \textbf{COrigami} \citep{corigami} has a \emph{citation missing} as well: it appears in
the metrics survey as the source of flat-foldability as a boolean and a VLM aesthetic score. Same
action needed.

\paragraph{Complexity and theory of flat folding.}
\citet{bern-hayes} is the citation under the paper's central claim: geometry does not determine
layer order, and deciding overlap order is NP hard even given a valid mountain-valley assignment.
\citet{arkin-map} defines simple foldability; the action space comes from here.
\citet{akitaya-infinite} covers the all-layers model, which corresponds to sheet-metal bending.
\citet{akitaya-flatfolder} is the Flat-Folder paper; deciding a global flat-folded state is NP
hard. See also \citet{mixed-orthogonal}, \citet{layer-algebra} and \citet{flat-folding-graphs}.
\citet{turing} is a defensive citation only, per Section~\ref{sec:objections}. \citet{continuous}
separates folded state from folding motion: reachability is free, so the hard question is the
discrete step structure; its predecessor is \citet{reaching-folded}.

\paragraph{The crease-pattern-to-sequence problem itself.}
\citet{akitaya-cp2seq} named the problem, by the group whose tools the field now uses: reflection
paths, graph rewriting, step graphs; a frog base explodes to 22,665 nodes and 30 minutes, and the
authors' own future work asks for the priority heuristic this paper's models are being asked to
supply. \citet{creasy} is an open-source implementation of the above and is usable as a symbolic
baseline. \texttt{[TO RUN if used]}

\paragraph{Counting folded states.}
\citet{oeis-a000136} lists the foldings of a strip of stamps: 1, 2, 6, 16, 50, 144, 462, 1392,
4536, 14060, \ldots{} There is no closed form. See also A001011 and A001416, and
\citet{lucas,koehler,lunnon,meanders}; meanders and stamp foldings are the same combinatorial
object.

\paragraph{Spatial reasoning context.}
\citet{spatial-survey} surveys the area. \citet{vot} names the mental-imagery question this
benchmark can turn into a measurement.

\paragraph{Tools, formats and data.}
\citet{flatfolder} decides and enumerates flat-folded states and is the source of the four
constraint classes cited in Section~\ref{sec:env-refusals}. \citet{fold-format} is the de facto
standard used by every artifact we ship. \citet{oripa} is a crease-pattern editor with folded-form
estimation. \citet{purelandfold} is not used as training or evaluation data here; it is cited as
the closest existing sequence-level corpus and as FoldingAgent's data.

\paragraph{Statistical physics of random flat-foldability.}
Optional support for the claim that multiple valid states is the normal case rather than a
curiosity, to be included only if Section~\ref{sec:problem-states} needs reinforcement:
\citet{spin-model} and \citet{random-locally}.

\paragraph{Deliberately not cited.}
Thick folding (Ku \& Demaine 2016), bar-and-hinge mechanics, MERLIN2, SWOMPS, Sim-FAST-PY, and the
creased-sheet mechanics literature all belong to the thickness and mechanics line of work. They are
a separate paper with a verifier several orders of magnitude more expensive, and citing them here
would invite a reviewer to ask why this paper does not do that one.

\section{Editorial rules for this draft}
\label{app:editorial}

Not part of the paper; carried over from \texttt{paper/DRAFT.md} so that nothing is lost in the
conversion. Each of these has been drafted wrongly at least once in this project.

\begin{enumerate}
\item \textbf{Never claim legal moves are rare.} The enumeration refutes it, per
  Section~\ref{sec:problem-task}.
\item \textbf{Never claim the model performs 3D reconstruction.} Say spatial reasoning; pose ``is a
  three-dimensional representation necessary or merely sufficient'' as a research question. Avoid
  homotopy and isotopy language entirely.
\item \textbf{Turing completeness is defensive only}, per Section~\ref{sec:objections}.
\item \textbf{Cite OrigamiSpace as an arXiv preprint} unless a venue is confirmed;
  \texttt{papers.md} recorded a venue that its arXiv record does not carry.
\item \textbf{Every number carries a \texttt{fact:} tag} and \texttt{doccheck.mjs} must pass before
  submission. The tags survive in this file as LaTeX comments beside each number.
\item \textbf{Sections \ref{sec:experiments} to \ref{sec:results} stay empty until the runs exist.}
  If they cannot be filled in time, the paper is submitted as a benchmark paper and the ablation
  sentence leaves the abstract.
\end{enumerate}

\section{Draft status}
\label{app:status}

Carried over from the header of \texttt{paper/DRAFT.md}. Sections \ref{sec:intro} to
\ref{sec:objections} and \ref{sec:limitations} to \ref{sec:reproducibility} are written. Sections
\ref{sec:experiments}, \ref{sec:ablations} and \ref{sec:results} are empty and marked
\texttt{[TO RUN]}, each saying which run would fill it. \textbf{Nothing goes into those three
sections that has not been run.}

Ownership: \texttt{paper/DRAFT.md} owns the paper's prose. It cites \texttt{DATASET.md},
\texttt{EXPERIMENTS\_SETUP.md} and \texttt{notes/plan/*} rather than restating them. Numbers carry
the repo's inline \texttt{fact:} tags, so \texttt{doccheck.mjs} guards them there as everywhere
else. Appendix~\ref{app:editorial} holds the editorial rules that are about writing rather than
content.

\end{document}
